% 分类目录：dl
\documentclass[12pt, a4paper]{article}
\usepackage[margin=2.5cm]{geometry}
\usepackage{ctex}
\usepackage{amsmath, amssymb}
\usepackage{listings}
\usepackage{xcolor}
\usepackage{tabularx}
\usepackage{tikz}
\usetikzlibrary{shapes, arrows.meta, positioning, calc}

% 定义样式
\tikzset{
    layer/.style={rectangle, draw, thick, minimum width=3cm, minimum height=0.8cm, align=center, rounded corners=2pt},
    attn/.style={layer, fill=blue!10},
    ffn/.style={layer, fill=orange!10},
    norm/.style={layer, fill=green!10},
    emb/.style={layer, fill=gray!20},
    pos/.style={circle, draw, thick, inner sep=2pt, fill=yellow!20},
    arrow/.style={-Stealth, thick}
}

\lstset{
    language=Python,
    basicstyle=\ttfamily\small,
    keywordstyle=\color{blue},
    commentstyle=\color{green!50!black},
    stringstyle=\color{red},
    showstringspaces=false,
    breaklines=true,
    numbers=left,
    numberstyle=\tiny\color{gray},
    frame=single
}

\title{知识点：Transformer}
\author{}
\date{}

\begin{document}
\maketitle

\section{核心定义}
\textbf{中文定义}：Transformer 是一种基于自注意力机制（Self-Attention Mechanism）的深度学习架构，摒弃了传统的循环神经网络（RNN）和卷积神经网络（CNN），通过全局注意力机制捕获序列数据中的长距离依赖关系。

\textbf{英文定义}：Transformer is a deep learning architecture based entirely on self-attention mechanisms, dispensing with traditional recurrence and convolutions entirely to capture long-range dependencies in sequential data.

\textbf{提出时间与作者}：2017年，由 Google 团队（Ashish Vaswani 等人）提出。

\textbf{原始关键论文}：\textit{Attention Is All You Need} (NeurIPS 2017)

\textbf{技术体系定位及历史演进}：
Transformer 最初作为自然语言处理（NLP）领域机器翻译任务的 Seq2Seq 模型提出，旨在解决 RNN 难以并行计算以及长跨度依赖（Long-term Dependency）容易遗忘的问题。由于其卓越的特征提取能力和极强的并行扩展性，Transformer 迅速取代了 LSTM/GRU，成为 NLP 的事实标准（催生了 BERT、GPT 等大模型）。随后，其应用领域扩展至计算机视觉（Vision Transformer, ViT）、多模态学习以及强化学习等各个 AI 领域，成为现代深度学习的基础架构。

\section{核心原理与底层逻辑}
Transformer 的核心在于\textbf{多头自注意力（Multi-Head Self-Attention）}机制，加上前馈神经网络（Feed-Forward Networks）、残差连接（Residual Connection）和层归一化（Layer Normalization）。

\subsection{数据流向简述}
假定输入序列长度为 $L$，嵌入维度为 $D$，批量大小为 $B$。
\begin{enumerate}
    \item \textbf{输入张量}：原始词元进入 Embedding 层，转化为形状为 $(B, L, D)$ 的张量，并加上相同形状的 \textbf{位置编码（Positional Encoding）}，以注入序列的顺序信息。
    \item \textbf{自注意力计算}：张量通过线性变换分别投影出 $Q$ (Query), $K$ (Key), $V$ (Value)，形状均为 $(B, L, D)$（多头情况下会被拆分成多份）。注意力机制计算词元之间的相关性得分并加权求和，输出形状依旧是 $(B, L, D)$。
    \item \textbf{前馈网络与输出}：经过残差连接与 LayerNorm 后，送入两层全连接网络（内部维度通常放大 4 倍），再经残差与 LayerNorm 最终输出形状为 $(B, L, D)$ 的张量。
\end{enumerate}

\subsection{算法执行步序}
\begin{enumerate}
    \item 将输入 Token 映射为词向量 $X$。
    \item 生成位置编码 $PE$，并计算 $X \leftarrow X + PE$。
    \item 将 $X$ 投影为三个矩阵：查询（$Q$）、键（$K$）、值（$V$）。
    \item 计算缩放点积注意力（Scaled Dot-Product Attention）：$Q$ 与 $K^T$ 相乘，除以放缩因子，经过 Softmax 得到注意力权重。
    \item 注意力权重与 $V$ 相乘，得到加权上下文向量。
    \item 加入残差连接并进行 LayerNorm：$\text{Output}_1 = \text{LayerNorm}(X + \text{Attention}(X))$。
    \item 输入到全连接前馈网络（FFN），再次加入残差连接并进行 LayerNorm：$\text{Output}_2 = \text{LayerNorm}(\text{Output}_1 + \text{FFN}(\text{Output}_1))$。
\end{enumerate}

\section{严谨数学公式与推导}
Transformer 的最核心公式是缩放点积注意力（Scaled Dot-Product Attention）：
$$
\text{Attention}(Q,K,V) = \text{softmax}\left(\frac{QK^T}{\sqrt{d_k}}\right)V
$$
\textbf{符号说明}：
\begin{itemize}
    \item $Q \in \mathbb{R}^{L \times d_k}$：Query 矩阵，表示当前词元的查询请求。
    \item $K \in \mathbb{R}^{L \times d_k}$：Key 矩阵，表示序列中所有词元的被查询标识。
    \item $QK^T \in \mathbb{R}^{L \times L}$：两者相乘得到注意力分数矩阵，表示序列中第 $i$ 个词元对第 $j$ 个词元的注意力大小。
    \item $\sqrt{d_k}$：缩放因子（Scaling Factor），$d_k$ 是 $K$ 的维度。作用是防止点积结果过大，导致 Softmax 函数进入梯度极小的饱和区。
    \item $V \in \mathbb{R}^{L \times d_v}$：Value 矩阵，表示词元的实际内容表示。
    \item \text{softmax}：对每一行（每个词元的注意力分布）进行归一化，使其和为 1。
\end{itemize}

多头注意力（Multi-Head Attention）则是将 $Q, K, V$ 映射到多个不同的线性子空间中分别计算注意力，再将结果拼接融合：
$$
\text{MultiHead}(Q,K,V) = \text{Concat}(\text{head}_1, \dots, \text{head}_h) W^O
$$
其中，$\text{head}_i = \text{Attention}(Q W_i^Q, K W_i^K, V W_i^V)$。

\section{模型结构示意图 (Full Architecture)}

\begin{figure}[h!]
\centering
\resizebox{0.88\textwidth}{!}{
\begin{tikzpicture}[node distance=1.8cm, auto]

    % --- Encoder ---
    \node [emb] (input_emb) {Input Embedding};
    \node [pos] (pos_enc_in) [right=0.2cm of input_emb] {+};
    \node [below=0.1cm of pos_enc_in, font=\tiny] {Positional Encoding};
    
    \node [attn, above=of input_emb] (multi_head_enc) {Multi-Head Self-Attention};
    \node [norm, above=0.6cm of multi_head_enc] (norm1_enc) {Add \& Norm};
    \node [ffn, above=of norm1_enc] (ffn_enc) {Feed Forward};
    \node [norm, above=0.6cm of ffn_enc] (norm2_enc) {Add \& Norm};
    
    \draw [arrow] (input_emb) -- (multi_head_enc);
    \draw [arrow] (multi_head_enc) -- (norm1_enc);
    \draw [arrow] (norm1_enc) -- (ffn_enc);
    \draw [arrow] (ffn_enc) -- (norm2_enc);
    
    % Residuals Encoder
    \draw [arrow] ($(input_emb.north)$) -- ++(-1.8,0) |- (norm1_enc.west);
    \draw [arrow] ($(norm1_enc.north)$) -- ++(-1.8,0) |- (norm2_enc.west);
    
    \node [left=1.5cm of norm2_enc, font=\small\bfseries] {Encoder ($\times N$)};

    % --- Decoder ---
    \node [emb, right=4.0cm of input_emb] (output_emb) {Output Embedding\\(Shifted Right)};
    \node [pos] (pos_enc_out) [right=0.2cm of output_emb] {+};
    
    \node [attn, above=of output_emb] (masked_attn) {Masked Multi-Head\\Self-Attention};
    \node [norm, above=0.6cm of masked_attn] (norm1_dec) {Add \& Norm};
    
    \node [attn, above=of norm1_dec] (cross_attn) {Multi-Head\\Cross-Attention};
    \node [norm, above=0.6cm of cross_attn] (norm2_dec) {Add \& Norm};
    
    \node [ffn, above=of norm2_dec] (ffn_dec) {Feed Forward};
    \node [norm, above=0.6cm of ffn_dec] (norm3_dec) {Add \& Norm};
    
    \node [layer, fill=purple!10, above=of norm3_dec] (linear) {Linear \& Softmax};
    
    \draw [arrow] (output_emb) -- (masked_attn);
    \draw [arrow] (masked_attn) -- (norm1_dec);
    \draw [arrow] (norm1_dec) -- (cross_attn);
    \draw [arrow] (cross_attn) -- (norm2_dec);
    \draw [arrow] (norm2_dec) -- (ffn_dec);
    \draw [arrow] (ffn_dec) -- (norm3_dec);
    \draw [arrow] (norm3_dec) -- (linear);
    
    % Residuals Decoder
    \draw [arrow] ($(output_emb.north)$) -- ++(1.8,0) |- (norm1_dec.east);
    \draw [arrow] ($(norm1_dec.north)$) -- ++(1.8,0) |- (norm2_dec.east);
    \draw [arrow] ($(norm2_dec.north)$) -- ++(1.8,0) |- (norm3_dec.east);
    
    % Cross-Attention Connection
    \path (norm2_enc.east) -- (cross_attn.west) coordinate[midway] (mid_point);
    \draw [arrow, dashed, ultra thick, blue!60] (norm2_enc.east) -- (mid_point |- norm2_enc.east) 
          node[above, font=\tiny, fill=white, inner sep=1pt] {Encoder Output (K, V)} 
          -| (cross_attn.west);
    
    \node [right=1.0cm of norm3_dec, font=\small\bfseries] {Decoder ($\times N$)};

\end{tikzpicture}}
\caption{Transformer Full Architecture (Attention Is All You Need)}
\end{figure}

\section{关键代码片段 (PyTorch实现)}
以下为 Transformer 中最核心的多头自注意力的简化实现（便于展示原理，省略了多头拆分的重塑操作细节）：

\begin{lstlisting}
import torch
import torch.nn as nn
import torch.nn.functional as F
import math

class SelfAttention(nn.Module):
    def __init__(self, d_model, heads):
        super().__init__()
        self.d_model = d_model
        self.heads = heads
        self.d_k = d_model // heads
        
        # 定义计算 Q, K, V 的线性层
        self.q_linear = nn.Linear(d_model, d_model)
        self.k_linear = nn.Linear(d_model, d_model)
        self.v_linear = nn.Linear(d_model, d_model)
        
        # 最后的输出映射层
        self.out = nn.Linear(d_model, d_model)
        
    def forward(self, x, mask=None):
        # x shape: (B, L, d_model)
        B, L, _ = x.shape
        
        # 1. 线性映射并拆分为多头 (B, L, heads, d_k) -> 调换维度 (B, heads, L, d_k)
        # 为了代码简洁，用伪操作表示 reshape 和 transpose
        Q = self.q_linear(x).view(B, L, self.heads, self.d_k).transpose(1, 2)
        K = self.k_linear(x).view(B, L, self.heads, self.d_k).transpose(1, 2)
        V = self.v_linear(x).view(B, L, self.heads, self.d_k).transpose(1, 2)
        
        # 2. 计算 Scaled Dot-Product Attention: Q * K^T / sqrt(d_k)
        # attention shape: (B, heads, L, L)
        scores = torch.matmul(Q, K.transpose(-2, -1)) / math.sqrt(self.d_k)
        
        if mask is not None:
            scores = scores.masked_fill(mask == 0, -1e9)
            
        # 注意力权重归一化
        attn_weights = F.softmax(scores, dim=-1)
        
        # 3. 乘以 Value
        # output shape: (B, heads, L, d_k)
        context = torch.matmul(attn_weights, V)
        
        # 4. 拼接多头并线性映射恢复原始维度
        context = context.transpose(1, 2).contiguous().view(B, L, self.d_model)
        output = self.out(context) # (B, L, d_model)
        
        return output
\end{lstlisting}

\section{深度面试题与解答}
\subsection{基础概念}
\textbf{问：Transformer 为什么要加入位置编码及其运作方式？}\\
答：因为自注意力机制（Q、K、V相乘）计算两两之间的关系时不关注它们在序列中的物理距离，这是一种完全等距的“集合”操作。如果不加入位置编码，模型会将“I like ML”和“ML like I”视作完全一致。原论文采用了正弦和余弦函数的绝对位置编码，通过将位置编码直接与输入词嵌入相加（$dim$一致）来给每个位置打上标签。

\subsection{进阶原理}
\textbf{问：为什么点积注意力在经过 Softmax 之前要除以 $\sqrt{d_k}$？}\\
答：当维度 $d_k$ 很大时，相互独立且服从标准正态分布的 $Q$ 和 $K$ 的点积，均值虽为 0，但方差会增大至 $d_k$。巨大的方差导致分布呈现两极分化，某些元素值极大。此时传入 Softmax 函数，会导致大的部分接近 1、其余接近 0，即 Softmax 函数陷入了梯度极小的平滑区域。除以 $\sqrt{d_k}$ 相当于将点积结果的方差缩放回 1，使得梯度流更加稳定。

\subsection{工程实战}
\textbf{问：在 Transformer 中 LayerNorm 放置的位置有哪几种？各有什么优劣？}\\
答：主要分为 Post-LN（原论文，在残差相加后执行 LN）和 Pre-LN（目前的绝对主流如 GPT-3，在 Attention 和 FFN 之前执行 LN）。\textbf{Post-LN} 对权重的初始分布极其敏感，梯度容易在靠近输入端消失，训练极不稳定，常需要设置复杂的 Warmup 学习率调度机制；\textbf{Pre-LN} 的梯度流动从后到前是一致的，不管层数多深，无需 Warmup 也可以直接稳定收敛。

\subsection{坑点分析}
\textbf{问：Transformer 处理长文本时面临的最大性能瓶颈是什么？如何解决？}\\
答：瓶颈在于计算自注意力的复杂度是 $O(L^2 \times D)$，空间复杂度（存储 Attention Map）也是 $O(L^2)$。当需要处理数十万 token 时显存会瞬间爆炸。解决方式有：
1. 稀疏注意力（如 Longformer 仅计算局部窗）。
2. 线性关联矩阵（如 Linear Transformer 重新排列向量乘积顺序）。
3. 状态复用或内存压缩机制（如 Transformer-XL）。
4. (近期绝对主流) 状态空间模型（Mamba）或是 FlashAttention 高效 IO 并发算法。

\section{优缺点与适用场景}

\begin{table}[h]
\centering
\begin{tabular}{|p{0.2\textwidth}|p{0.75\textwidth}|}
\hline
\textbf{维度} & \textbf{分析} \\
\hline
\textbf{优势} & 1. 全局视野：一步直接建立任意长距离依赖关系。\\
& 2. 高度并行化：训练时不需要像 RNN 那样等待前序步骤完成，硬件利用率极高。\\
& 3. 泛化与扩展性：模型容量扩展到千亿参数也未见明显天花板（Scaling Law 表观有效）。\\
\hline
\textbf{局限} & 1. 计算复杂度过高：序列长度的二次方复杂度阻碍了极致长文应用。\\
& 2. 缺乏归纳偏置（Inductive Bias）：相比CNN对图像平移的不变性假设，Transformer 较难在小数据集上收敛，极为依赖海量数据的暴力压制。\\
\hline
\textbf{最佳适用场景} & 海量数据下的大模型预训练（NLP 的语言模型、代码生成）、机器翻译、高质量文本生成；近期的图像生成的 DiT 架构。\\
\hline
\textbf{绝对慎用场景} & 1. 训练数据极少且无预训练支持的孤立小任务。\\
& 2. MCU（单片机）等边缘侧算力极度受限设备的实时推理。\\
\hline
\end{tabular}
\end{table}

\section{参考文献与拓展}
\begin{enumerate}
    \item \textbf{核心经典论文}: Vaswani, A., et al. (2017). \textit{Attention is all you need}. NeurIPS. 
    \item \textbf{代码实现参考}: \texttt{torch.nn.MultiheadAttention} 和 \texttt{torch.nn.TransformerEncoderLayer} 官方源码。
    \item \textbf{必读进阶}: \textit{The Illustrated Transformer}（Jay Alammar，目前最生动的图文详解）。
    \item \textbf{变体探索}: FlashAttention (Dao et al., 2022) 解决硬件 IO 通信瓶颈的终极方案。
\end{enumerate}

\end{document}
